\section{Core Technologies}

In addition to investigating new techniques for managing
system resources and data, we will attempt to share a common image
representation for both resource and data management.
Traditionally, the image processing community and the database 
community have been developing their own image representations, i.e.,
the image processing community proposes 
compression algorithms and the database community proposes image
feature representations (e.g., color histograms).
A fascinating research direction may be to build search indexes 
on the compressed data directly. We are especially interested
in exploring {\em layered coding} techniques that can support
both scalable content degradation and scalable content-based indexing.

Finally, our design goal is not only to provide the 
interactive functions, but also
provide these functions with minimum resource requirement
(hence at minimum cost) and with minimum initial latency.
In the following sections, I enumerate the technologies 
we plan to work on, beginning with our efforts in resource
management. 

\subsubsection{Indexing for Temporally Related Data}

When a data object
is referenced, it and all its temporally related objects should be
found quickly, and once requested, be retrieved efficiently. 
Indexing temporally related objects can be tricky.
In principle, all temporally related object should be 
stored contiguously on disk to minimize the disk latency 
necessary to retrieve them.
However, since an object is often shared in many scenes,
replication may not always be desirable. One alternative is
storing shared objects in the middle tracks so their
access latency is constrainted. Another alternative
is storing shared objects on the outer tracks to take
advantage of their higher data transfer rates. 
We will experiment with various data placement policies to
examine if some rules of thumb can work most of the time.

It seems that building indexes at the scene level is
appropriate since a scene is a reasonable search 
unit (why this is so is discussed further 
in the content-based retrieval section,
which follows). Once a scene can be located, the 
scene description can tell where its data objects
reside. I believe that the combination of hashing and XML
may be an appropriate structure for indexing scenes.
A hash table provides fast scene lookup; XML
provides the flexibility to describe semistructured
objects in a scene. We will investigate how to
deal with efficient updates (primarily necessary because of changes
of the object's location)
and queries on XML structures \cite{Abi97,Bun97,MW99}.
Another challenging problem and a new research direction
is performing similarity queries on semi-structured 
data that is in XML format.

One useful feature that the receiver can support is video
editing. A video editing tool can combine scenes from
different videos into a customized clip. Supporting
video editing introduces additional scene management issues.
When a scene is shared by more than one clip, the scene
should not be copied entirely. A scene is copied only 
when its substantial content is modified. Otherwise,
we may want to create a shell scene, which maintains
links to its parent scene and keeps only those objects that it
modified. 

\subsubsection{Indexing for Content-Based Retrieval}

To support effective time-shift operations and general
multimedia database queries, we
intent to build a content-based video scene
retrieval system. 
A video can be divided into a number of sequences,
a sequence a number of scenes, and a scene a number
of frames.  We believe that a scene should be a reasonable
query unit. Supporting queries on the frame level 
is simply too expensive (a one-hour video
at 25 frames/second frame rate has $90$K frames).
Furthermore, indexing on the frame level 
cannot describe motion, e.g., a bus is approaching.
Supporting queries on the sequence level
may be too coarse and may hurt query precision.
Therefore, we plan to build scene descriptors
and build indexes on these scene descriptors.

Building such a system faces
four challenging steps: 1) selecting and extracting features from
multimedia data objects, 2) clustering the extracted features,
3) indexing objects based on their feature clusters, and 4) devising
query processing steps that find matches efficiently
in the database.

Selecting and extracting features can be complicated.
Fortunately, much low-hanging fruit in this area is ready to be plucked.
We plan to rely
on existing tools as much as possible. For feature clustering,
a large number of papers \cite{img11,img12,img13}
have been published with many of
ideas. Instead of coming up with yet another feature
cluster, we will experiment with a variance control model
that will decide how fine the features must be clusters to attain
good similarity results.
Take image searching as an example.
If one wants to find near-copies of an image,
the feature cluster should record finer details 
of the color distribution on the query image. If one wants to find
images of similar colors that are not necessarily
at the same spots, the feature vector can record 
coarser color distribution. 
Extensive exploration and experiments will 
be performed to verify if this variance controlled 
approach can be fruitful.

Content-based retrieval 
is inherently expensive, especially when
it needs a large number of dimensions to characterize the
target object (e.g., an image or a video scene).  
There is simply no way to
build an index on disk such that {\em all} nearest neighbors to {\em
any} query point are physically near on disk.  
Fortunately, in many cases it is
sufficient to perform an {\em approximate search} that returns many
but not all nearest neighbors~\cite{AM94,Cla94,IM98,Jk97,KOR98}.  (A
feature vector is often an approximate characterization of an object,
so we are already dealing with approximations anyway.)  
I have recently proposed a
clustering/indexing combined scheme that achieves approximate
similarity search with high efficiency \cite{Lch99}.
This proposed clustering/indexing combined scheme should
provide a solid foundation for continued work,
but many challenges remain. We are, for one, not
convinced that the Euclidean distance
is the only way to measure similarity for multimedia data.
Thus, we are investigating other similarity measures
(e.g., Hamming distance, etc.). For another,
we think similarity should be personalizable.
We will thus investigate new measures for similarity.

Finally, 
for query processing, we will investigate 
data models, query languages and query optimization techniques.

%\subsection{Image Coding and Packetization}
%\label{sec-image}
%
%The researchers in the compression
%camp have been devised compression techniques 
%for coding efficiency and robustness. The
%researchers in the database camp have been devised
%indexing data organization schemes for efficient data retrieval.
%Researchers in two camps often characterize an image 
%differently. Consequently, an image may have one
%representation (e.g., JPEG) in transmission and
%another representation (e.g., color histograms)
%for the indexing phase. 
%It would be nice if two efforts can
%share one common data representation. At this time, wavelets
%seem to be an attractive choice.
%Wavelets' scalability fits well
%for supporting both scalable content degradation and
%scalable content-based indexing.
%
%%While studying image representation, we also want to
%investigate how an image should be packetized to
%be more error resilient.
%One idea is that for the base image/audio/video layer, 
%forward error correction (FEC) code can be applied to
%improve its error resiliency. For progressive layers, 
%however, we may not want to
%spend additional bits on channel encoding. Instead, we can rely on 
%error concealment and reconstruction techniques to
%recover an image after encountering channel errors.
%
%We will also investigate image-based rendering techniques.
%At this end, our focus is on video panoramas for supporting
%spatial viewing interactivity.
%For a single camera, supporting
%video panoramas requires a method for locating
%frames that contain smooth
%camera pans or tilts, from which the panoramic
%view can be synthesized.
%

%\subsection{Putting All the Pieces Together}


